% Capitulo: Fundamentos que não envelhecem
\chapter{Fundamentos que não envelhecem}
\label{cap:fundamentos}

Ferramentas envelhecem rápido. A ferramenta de BI que alguém dominava há dez anos
provavelmente já foi substituída, e o assistente de IA de hoje vai parecer antigo em dois
anos. Algumas habilidades, porém, atravessam todas as ondas que a Parte~\ref{parte:origens}
descreveu sem perder valor. Este capítulo reúne cinco delas: modelagem, estatística e
pensamento crítico, qualidade de dados, comunicação e conhecimento do negócio. Nenhuma
depende de produto, e todas ficaram mais importantes com a IA, porque são justamente o que
permite avaliar o que a máquina produz.

\section{Modelagem}
\label{sec:fund-modelagem}

Modelar dados é decidir como o mundo vai ser representado em tabelas: o que é uma
entidade, o que identifica cada registro, como as coisas se relacionam, o que uma linha
representa. O capítulo~\ref{cap:dw} apresentou duas formas clássicas de fazer isso, a
normalizada e a dimensional. Aqui interessa o que elas têm em comum.

Toda boa modelagem responde a três perguntas antes de qualquer outra. Qual é o grão, ou
seja, o que exatamente uma linha representa? Qual é a chave, ou seja, o que torna cada
linha única? E como esta tabela se liga às outras? Quem pergunta isso por hábito evita a
maior parte dos erros de análise, inclusive os que a IA comete com toda a confiança, como o
frete multiplicado do capítulo anterior (\veja{sec:consulta-errada}).

Modelagem é também a habilidade que mais se valoriza com a IA. Quando escrever código fica
barato, o que diferencia um projeto bom de um ruim é o desenho: as tabelas certas, com o
grão certo e os nomes certos. A Parte~\ref{parte:modelagem} mostra que esse desenho muda
conforme o uso, mas as três perguntas continuam as mesmas.

\section{Estatística e pensamento crítico}
\label{sec:fund-estatistica}

Não é preciso ser estatístico para trabalhar com dados, mas é preciso pensar como um em
alguns momentos. Na maior parte do tempo, isso se resume a poucas ideias, aplicadas com
disciplina.

A primeira é olhar a distribuição antes do resumo. Uma média esconde tudo o que importa
quando os dados são assimétricos: o tempo médio de entrega pode estar ótimo enquanto um
em cada dez pedidos atrasa uma semana. Mediana, percentis e um histograma simples contam a
história que a média esconde.

A segunda é sempre comparar. Um número sozinho não diz nada. Cresceu em relação a quê? Está
alto em comparação com o quê? A variação está dentro do normal ou é sinal de alguma
mudança? Boa parte do trabalho de análise é escolher a comparação certa.

A terceira é desconfiar de relações simples. Correlação não é causalidade, e a frase é tão
repetida que perdeu a força. Um exemplo clássico ajuda a recuperá-la. Em 1973, os dados de
admissão na pós-graduação da Universidade da Califórnia em Berkeley mostravam que homens
eram aceitos em proporção maior do que mulheres. Olhando departamento por departamento, o
viés desaparecia, e em vários deles as mulheres tinham taxa de aceitação um pouco maior. A
explicação: as mulheres se candidatavam mais aos departamentos mais concorridos, que
recusavam a maioria dos candidatos de qualquer sexo \parencite{bickel1975}. O agregado
contava uma história; os grupos, outra. Esse fenômeno, em que uma tendência se inverte
quando os dados são separados em grupos, é conhecido como paradoxo de Simpson.

Num painel de vendas, o mesmo efeito aparece quando a margem média cai porque a empresa
vendeu mais de um produto de margem baixa, mesmo que a margem de cada produto tenha subido.
Quem não pensa em composição tira a conclusão errada de um número correto.

A quarta ideia é lembrar que números podem ser apresentados de forma enganosa, por
intenção ou por descuido. O pequeno livro de \textcite{huff1954}, com mais de setenta anos,
continua sendo a melhor introdução ao assunto: eixos cortados, médias escolhidas a dedo,
amostras viciadas. E vale conhecer os atalhos do nosso próprio raciocínio, que tendem a
construir uma história coerente com o pouco que está à vista e a ignorar o que não está
\parencite{kahneman2011}. Um assistente de IA que responde com segurança e fluência explora
exatamente essa tendência.

\section{Qualidade de dados}
\label{sec:fund-qualidade}

Qualidade de dados costuma ser tratada como exatidão: o valor está certo ou errado. Um
estudo clássico mostrou que, para quem usa os dados, qualidade é bem mais do que isso.
\textcite{wang1996} perguntaram a consumidores de dados o que tornava um dado bom e
organizaram as respostas em quatro grupos: qualidade intrínseca (o dado é correto e
confiável), qualidade contextual (é relevante, atual e completo para a tarefa), qualidade
de representação (é compreensível e consistente) e acessibilidade (dá para obter, com
segurança).

A consequência prática é que qualidade depende do uso. Um dado atualizado uma vez por mês
tem ótima qualidade para o fechamento mensal e péssima para o acompanhamento diário. Um
cadastro com o nome do cliente escrito de três formas é aceitável para contar pedidos e
inaceitável para somar a receita por cliente. Perguntar \enquote{este dado é bom?} sem
perguntar \enquote{bom para quê?} não leva a lugar nenhum.

Na prática do dia a dia, isso vira hábitos simples: conferir se as chaves são únicas,
contar nulos, comparar totais com a fonte, olhar a data da última atualização antes de
confiar num número. Esses hábitos viram testes automáticos no capítulo~\ref{cap:testes} e
uma disciplina de governança no capítulo~\ref{cap:governanca}.

\section{Comunicação}
\label{sec:fund-comunicacao}

Uma análise que ninguém entende não muda nenhuma decisão. Comunicar é parte do trabalho, e
não um enfeite no fim dele.

A recomendação mais útil sobre escrita para quem trabalha com dados vem da consultoria:
comece pela conclusão. Barbara Minto propôs organizar qualquer texto como uma pirâmide,
com a ideia principal no topo e os argumentos que a sustentam embaixo
\parencite{minto2009}. Um e-mail que começa com \enquote{as vendas caíram 8\% por causa da
ruptura de estoque em duas regiões} e só depois mostra os números é lido; um que começa
pela metodologia e chega à conclusão no quinto parágrafo, não.

O mesmo vale para gráficos e apresentações (\veja{sec:visualizacao}): saber quem vai ler,
o que essa pessoa precisa decidir, e colocar a mensagem no título, não numa legenda que
ninguém lê. E vale, cada vez mais, para a forma como se escreve para a IA. Um pedido claro,
com contexto, objetivo e critério de qualidade, é um pequeno exercício de comunicação, e
quem escreve bem para pessoas costuma escrever bem para máquinas.

\section{Conhecer o negócio}
\label{sec:fund-negocio}

O quinto fundamento é o que menos se ensina e o que mais diferencia: entender o negócio em
que os dados nascem. Saber como a empresa ganha dinheiro, quais decisões dependem dos
números e onde a operação costuma falhar transforma uma análise correta numa análise útil.

Os quatro fundamentos anteriores ficaram mais acessíveis com a IA, que ajuda a modelar, a
calcular, a conferir e a escrever. Este é o único em que ela pouco ajuda, porque o
conhecimento de negócio está nas pessoas e nos processos da empresa, e não nos dados de
treinamento de um modelo. Por isso ele ganha o próximo capítulo inteiro.

\begin{leituras}
  \leitura{huff1954}{Curto, divertido e ainda atual. A melhor vacina contra gráficos e
    números enganosos.}
  \leitura{wang1996}{O estudo que mostrou que qualidade de dados, para quem usa, vai muito
    além da exatidão.}
  \leitura{minto2009}{O método da pirâmide para organizar textos e apresentações. Muda a
    forma de escrever qualquer e-mail de análise.}
  \leitura{kahneman2011}{Para entender os atalhos do raciocínio que levam a conclusões
    erradas, inclusive diante de números corretos.}
\end{leituras}
